\documentclass[12pt]{article}

% Layout.
\usepackage[top=1in, bottom=0.75in, left=1in, right=1in, headheight=1in, headsep=6pt]{geometry}

% Fonts.
\usepackage{mathptmx}
\usepackage[scaled=0.86]{helvet}
\renewcommand{\emph}[1]{\textsf{\textbf{#1}}}
\usepackage{graphicx}
\usepackage{wrapfig}

% TiKZ.
\usepackage{tikz, pgfplots}
\usetikzlibrary{calc}
\pgfplotsset{compat = newest}
\pgfplotsset{my style/.append style={axis x line=middle, axis y line=
middle, xlabel={$x$}, ylabel={$y$}, axis equal }}

% Misc packages.
\usepackage{amsmath,amssymb,latexsym}
\usepackage{graphicx}
\usepackage{array}
\usepackage{xcolor}
\usepackage{multicol}
\usepackage{enumerate}

% Commands to set various header/footer components.
\makeatletter
\def\doctitle#1{\gdef\@doctitle{#1}}
\doctitle{Use {\tt\textbackslash doctitle\{MY LABEL\}}.}
\def\docdate#1{\gdef\@docdate{#1}}
\docdate{Use {\tt\textbackslash docdate\{MY DATE\}}.}
\def\doccourse#1{\gdef\@doccourse{#1}}
\let\@doccourse\@empty
\def\docscoring#1{\gdef\@docscoring{#1}}
\let\@docscoring\@empty
\def\docversion#1{\gdef\@docversion{#1}}
\let\@docversion\@empty
\makeatother

% Headers and footers layout.
\makeatletter
\usepackage{fancyhdr}
\pagestyle{fancy}
\fancyhf{}
\lhead{\baselineskip 30pt
\emph{\@docdate\hfill\@doctitle}
\ifnum \value{page} > 1\relax\else\\
\emph{Name: \rule{3.5in}{1pt}\hfill \@docscoring}\fi}
\rfoot{\emph{\@docversion}}
\lfoot{\emph{\@doccourse}}
\cfoot{\emph{\thepage}}
\renewcommand{\headrulewidth}{0pt}
\makeatother

% Paragraph spacing
\parindent 0pt
\parskip 6pt plus 1pt

% Problem counters
\newcounter{probcount}
\newcounter{subprobcount}
\setcounter{probcount}{0}

\newcommand{\problem}[1]{%
\par\addvspace{4pt}%
\setcounter{subprobcount}{0}%
\stepcounter{probcount}%
\makebox[0pt][r]{\emph{\arabic{probcount}.}\hskip1ex}\emph{[#1 points]}\hskip1ex}

\newenvironment{subproblems}{%
\begin{enumerate}%
\setcounter{enumi}{\value{subprobcount}}%
\renewcommand{\theenumi}{\emph{\alph{enumi}}}}%
{\setcounter{subprobcount}{\value{enumi}}\end{enumerate}}

% Blanks
\newcommand{\blank}[1]{\rule{#1}{0.5pt}}
\newcommand{\mblank}[1]{\underline{\hspace{#1}}}

% Display helpers
\renewcommand{\d}{\displaystyle}
\newcommand{\ds}{\displaystyle}

% Document info
\doctitle{Math F251X: Quiz 7}
\docdate{March 19, 2026}
\doccourse{UAF Calculus I}
\docversion{Spring 2026}
\docscoring{\blank{0.8in} / 25}

\begin{document}

\emph{Please circle your instructor's name:}
\hfill James Gossell \hfill Gordon Williams \\

There are 25 points possible on this quiz.
Any outside materials are not allowed.
\emph{For full credit, show all work clearly.}

%------------------------------------------------
\problem{9} A 6-ft-tall person walks toward a wall at a rate of 3 ft/sec. A spotlight is located on the ground 24 ft from the wall. How fast does the height of the person’s shadow on the wall change when the person is 12 ft from the wall?

\hfill\includegraphics[width=3in]{image.png}
\newpage
\problem{8} Suppose I want to estimate the value of $3e^{0.1}$. 
\begin{subproblems}
\item Determine an appropriate $f(x)$ and $a$, and evaluate $L(x)=f(a)+f'(a)(x-a)$ at an appropriate point to estimate the value of $3e^{0.1}$. 
\vfill
\item Write down what you would need to enter into a calculator to determine the numerical error in the linear approximation you've just obtained.
\vspace{2cm}
\end{subproblems}
\problem{8} Consider the function \(f(x)=\dfrac{2x}{x^{2}+9}\).
\begin{subproblems}
\item Find the critical points of \(f(x)\)
 over the domain $[-1,5]$. 
Note that $f'(x)=\dfrac{18-2x^{2}}{\left(x^2+9\right)^2}$.
\vfill
\item Find the local and/or absolute extrema for the function $f(x)$ over the domain $[-1,5]$, or explain why none exist. Clearly label any you identify.
\vfill
\end{subproblems}
%\problem{Bonus: 2} 
%\vfill
%\vfill

%
%\begin{subproblems}
%\begin{multicols}{3}
%\item $\d \lim_{x \to -4^-} f(x) = \mblank{.5in}$
%\item $\d \lim_{x \to -4^+} f(x) = \mblank{.5in}$
%\item $\d \lim_{x \to -4} f(x) = \mblank{.5in}$
%\end{multicols}
%
%\begin{multicols}{3}
%\item $\d \lim_{x \to 2^+} f(x) = \mblank{.5in}$
%\item $\d f(4) = \mblank{.5in}$
%\item $\d \lim_{x \to 4} f(x) = \mblank{.5in}$
%\end{multicols}



\end{document}